\chapter{Structured Streaming Foundations}
\index{Structured Streaming}\index{micro-batch}\index{streaming source}\index{streaming sink}%
\index{output mode}\index{trigger}\index{availableNow}\index{checkpoint}\index{exactly-once}%
\begin{objectives}
\begin{itemize}
    \item Explain the infinite-table model and the micro-batch execution engine beneath it.
    \item Read from streaming sources (files, rate, Kafka preview) and write to Delta with checkpoints.
    \item Choose output modes (append, update, complete) and triggers (processing time, availableNow, continuous) deliberately.
    \item Operate streams: restart semantics, monitoring with listeners and metrics, and backpressure.
    \item Prove a stream survives kill-and-restart without loss or duplication.
    \item Design a sub-minute ingestion layer for a high-rate event source.
\end{itemize}
\end{objectives}

\begin{crosscloud}
Structured Streaming is Spark's \textbf{streaming engine}; the same model---micro-batches or
operators over an unbounded log---exists on every platform.

\vspace{2mm}
{\footnotesize
\begin{tabular}{@{}p{3.0cm}p{3.4cm}p{3.2cm}p{3.2cm}@{}}
\toprule
\textbf{Capability} & \textbf{Databricks (Spark)} & \textbf{AWS} & \textbf{Azure / GCP / Fabric} \\
\midrule
Engine & Structured Streaming & Flink (Managed Service) / Lambda & Stream Analytics / Dataflow (Beam) / Eventstream+Activator \\
Execution & Micro-batch (continuous optional) & True streaming (Flink) & Streaming / unified batch+stream (Beam) \\
Exactly-once & Checkpoint + idempotent sink & Flink checkpoints + 2PC sinks & At-least-once + dedup / exactly-once (Dataflow) \\
Replay & Source offsets in checkpoint & Kinesis iterators / Kafka offsets & Event Hubs checkpoints / Pub/Sub seek \\
Batch unification & \texttt{availableNow} trigger & Flink batch mode & Dataflow's core thesis \\
Monitoring & Stream metrics + listener & CloudWatch / Flink metrics & Azure Monitor / Dataflow UI / Fabric metrics \\
\bottomrule
\end{tabular}}

\vspace{1mm}
\textbf{Snowflake} approaches streaming with Snowpipe Streaming and dynamic tables (micro-batch
by design); \textbf{MongoDB} with change streams feeding processors.
\end{crosscloud}

\section{Week 13 Roadmap and Mental Model}
Part III processed data that had already arrived; Part IV processes data \emph{as} it arrives. Structured Streaming's achievement is making that feel like batch: the same DataFrame API, the same Catalyst optimizer, the same Delta sinks---over an input that never ends \cite{sparkStructuredStreaming}. Week 13 builds the execution model and the operational contract; Week 14 adds state and event time.

The mental model: \textbf{the stream is a table that grows; every trigger runs a small batch over the new rows}. Your query is planned once and executed incrementally; results land in the sink; the checkpoint remembers where the source was read to and what intermediate state exists. Understand that loop---trigger, plan, process new rows, commit offsets and state---and every streaming behavior becomes predictable.

\begin{definitionbox}
\textbf{Structured Streaming} treats a stream as an unbounded table. A query over it runs as a sequence of \textbf{micro-batches}: each trigger reads new source data since the last committed offset, applies the incrementally-executed query, writes the sink, and commits progress (offsets + state) to the \textbf{checkpoint}.
\end{definitionbox}

\begin{keyterms}
\textbf{Output mode}: what the sink receives---\texttt{append} (new rows only), \texttt{update} (changed rows), \texttt{complete} (whole result). \textbf{Trigger}: when micro-batches run---processing-time interval, \texttt{availableNow} (run once over backlog, then stop), or continuous. \textbf{Checkpoint}: the fault-tolerance state directory. \textbf{Backpressure}: bounding batch size so processing keeps up with input.
\end{keyterms}

\section{Monday: The Micro-Batch Model}
\subsection{Sources, Incremental Execution, Sinks}
A source (files via Auto Loader, Kafka, rate generator, Delta change feed) presents offsets; the engine reads a batch between the last committed offset and the current end; the query---planned once, including its stateful operators---processes incrementally; the sink (Delta in this program) commits transactionally. Because the sink write and the offset commit are coordinated through the checkpoint, a crash between them replays the batch---which is why sink idempotence (Chapter 2's MERGE, Delta's transactional replace) is the second half of exactly-once.

\subsection{Micro-Batch versus Continuous}
Micro-batch (the default) processes discrete chunks with seconds-to-minutes latency and full exactly-once support. \textbf{Continuous processing} runs long-lived operators for millisecond latency with at-least-once guarantees and a restricted operator set. The default answer in this program is micro-batch with a sane trigger; continuous is a deliberate, measured choice.

\begin{notebox}
Real-time does not mean continuous. Most ``real-time'' business requirements---dashboards refreshing every minute, alerts within a minute---are micro-batch requirements. Reserve continuous processing for measured sub-second mandates.
\end{notebox}

\section{Tuesday: Output Modes and Triggers}
\subsection{Choosing the Output Mode}
The mode is dictated by the query's semantics: a stateless filter \emph{appends}; an aggregation that keeps updating a key's value must \texttt{update} (or \texttt{complete}, rewriting the whole result---viable only for small results); aggregations over event time with watermarks append after the watermark closes the window (Chapter 14). Delta as a sink accepts append natively; update-style results use \texttt{foreachBatch} + MERGE (Chapter 16).

\subsection{Choosing the Trigger}
\texttt{processingTime="30 seconds"} for continuous low-latency flow; \texttt{availableNow=True} for scheduled, batch-priced streaming (process the backlog, stop, bill nothing until the next run---this is what makes Chapter 11's triggered ingestion economical); \texttt{Trigger.Once} for one-shot backfills. The trigger is a cost-latency dial you set per pipeline, per environment.

\begin{lstlisting}[language=Python,caption={The streaming skeleton used across Part IV.},label={lst:ch13-skeleton}]
stream = (spark.readStream.format("cloudFiles")
    .option("cloudFiles.format", "json")
    .option("cloudFiles.schemaLocation",
            "/Volumes/de_training/bronze/_schemas/clicks")
    .load("/Volumes/de_training/bronze/landing/clicks/"))

query = (stream.writeStream
    .format("delta")
    .outputMode("append")
    .option("checkpointLocation", "/Volumes/de_training/bronze/_ckpt/clicks")
    .trigger(availableNow=True)            # or: processingTime="30 seconds"
    .toTable("de_training.bronze.clicks"))

query.awaitTermination()
print(query.lastProgress)                  # input vs processing rate, offsets
\end{lstlisting}

\section{Wednesday: Checkpoints, Recovery, and Monitoring}
\subsection{What the Checkpoint Holds}
The checkpoint directory stores: source offsets (where to resume), state-store snapshots (for stateful operators, Chapter 14), and the query's metadata. Its rules, which this book treats as law: one checkpoint per query definition; never delete it casually (Chapter 11's pitfall); put it on governed storage; treat incompatible query changes as requiring a new checkpoint.

\subsection{Monitoring a Stream}
\texttt{query.lastProgress} and the streaming UI expose input rate, processing rate, batch duration, and offsets. \textbf{Falling behind} shows as input rate durably above processing rate; \textbf{unhealthy batches} show as rising duration. The two levers: more parallelism (partitions, executors) or smaller batches (\texttt{maxFilesPerTrigger}, \texttt{maxOffsetsPerTrigger}).

\begin{pitfall}
\textbf{Changing a stateful query in place.} Adding an aggregation key or changing window logic invalidates checkpointed state; the restart fails loudly (by design) or, worse with forced config, silently mis-aggregates. Version the query, new checkpoint, and a planned backfill---the same discipline as Chapter 11's checkpoint governance.
\end{pitfall}

\section{Thursday Hands-on Project: Kill It, Restart It, Prove It}
Build a file-sourced streaming pipeline, run it under three triggers, then kill it mid-stream and prove---from counts, not hope---that nothing was lost or duplicated.
\index{hands-on lab}\index{checkpoint!lab}\index{exactly-once!lab}

\begin{lstlisting}[language=Python,caption={Week 13 lab: restart-safety proof.},label={lst:ch13-lab}]
import uuid, time
from pyspark.sql import functions as F

# Producer: drop numbered JSON files into the landing path while streams run
def drop_batch(n, start):
    rows = [{"event_id": str(uuid.uuid4()), "seq": i,
             "ts": time.time()} for i in range(start, start + n)]
    (spark.createDataFrame(rows).write.mode("append")
     .json("/Volumes/de_training/bronze/landing/clicks/"))

drop_batch(500, 0)

q = (spark.readStream.format("cloudFiles")
     .option("cloudFiles.format", "json")
     .option("cloudFiles.schemaLocation", "/Volumes/de_training/bronze/_schemas/clicks")
     .load("/Volumes/de_training/bronze/landing/clicks/")
     .writeStream.format("delta")
     .option("checkpointLocation", "/Volumes/de_training/bronze/_ckpt/clicks")
     .trigger(processingTime="5 seconds")
     .toTable("de_training.bronze.clicks"))

drop_batch(500, 500)          # more data while running
q.stop()                      # the "crash" (mid-stream kill)
drop_batch(500, 1000)         # data arriving while the stream is DOWN

# Restart with the SAME checkpoint: must resume, not restart from zero
q2 = (spark.readStream.format("cloudFiles")
      .option("cloudFiles.format", "json")
      .option("cloudFiles.schemaLocation", "/Volumes/de_training/bronze/_schemas/clicks")
      .load("/Volumes/de_training/bronze/landing/clicks/")
      .writeStream.format("delta")
      .option("checkpointLocation", "/Volumes/de_training/bronze/_ckpt/clicks")
      .trigger(availableNow=True)
      .toTable("de_training.bronze.clicks"))

n = spark.table("de_training.bronze.clicks").count()
d = spark.table("de_training.bronze.clicks").dropDuplicates(["event_id"]).count()
assert n == 1500 == d, f"expected 1500 unique, got n={n} d={d}"
\end{lstlisting}

\subsection*{Deliverables}
\begin{itemize}
    \item The streaming pipeline with a governed checkpoint location.
    \item The kill-and-restart log: counts before stop, during downtime arrivals, after restart.
    \item The uniqueness assertion output proving 1500 rows, zero duplicates, zero loss.
\end{itemize}

\subsection*{Acceptance Criteria}
\begin{itemize}
    \item Restart from the same checkpoint resumes---it does not reprocess committed files.
    \item Final table equals exactly the produced events; dedup count equals raw count.
    \item \texttt{lastProgress} metrics (input rate, batch duration) captured for at least 5 batches.
    \item Trigger comparison documented: processing-time vs.\ \texttt{availableNow} on the same backlog.
\end{itemize}

\subsection*{Stretch Goals}
\begin{itemize}
    \item Add a \texttt{StreamingQueryListener} that logs batch metrics to a Delta table.
    \item Force a schema-incompatible query change; capture the restart failure and document the new-checkpoint procedure.
    \item Point the same source at a second sink with its own checkpoint; explain why that is safe while sharing a checkpoint is not.
\end{itemize}

\section{Friday Use Case: Sub-Minute Ingestion for Ad-Tech Events}
An ad-tech platform receives 50k events/second from bidder fleets. Downstream dashboards must see data within one minute; finance reconciles hourly; the events feed both real-time pacing and nightly billing.
\index{ad-tech}\index{streaming!design}

\subsection{Requirements and Assumptions}
Events arrive as JSON through a managed Kafka (3 partitions per topic is the current---wrong---configuration). Peak is 3x average. One minute p95 end-to-end latency. Nightly billing runs from the same Bronze table.

\begin{figure}[H]
    \centering
    \resizebox{\textwidth}{!}{%
    \begin{tikzpicture}[
        box/.style={rectangle, rounded corners, draw=primary, thick, fill=primary!7, align=center, minimum width=2.9cm, minimum height=0.95cm, font=\small\bfseries},
        store/.style={cylinder, draw=primary, shape border rotate=90, aspect=0.25, fill=primary!10, align=center, minimum width=2.3cm, minimum height=1.1cm, font=\small\bfseries},
        flow/.style={-{Latex[length=2mm]}, draw=primary, thick}
    ]
        \node[box] (bidders) at (0,0) {Bidder fleet\\50k events/s};
        \node[box] (kafka) at (3.6,0) {Kafka\\96 partitions};
        \node[box] (stream) at (7.2,0) {Structured Streaming\\30-s micro-batches\\capped offsets};
        \node[store] (bronze) at (11.2,0) {Bronze Delta\\(liquid clustered)};
        \node[box] (rt) at (14.6,0.8) {Pacing dashboards\\(1-min refresh)};
        \node[box] (bill) at (14.6,-0.8) {Nightly billing\\(batch, same table)};
        \draw[flow] (bidders) -- (kafka);
        \draw[flow] (kafka) -- (stream);
        \draw[flow] (stream) -- (bronze);
        \draw[flow] (bronze) -- (rt);
        \draw[flow] (bronze) -- (bill);
    \end{tikzpicture}%
    }
    \caption{Sub-minute ingestion: partition-scaled Kafka, capped micro-batches, one Bronze table serving both latencies.}
    \label{fig:ch13-adtech}
\end{figure}

\subsection{Architecture Decisions}
\textbf{Parallelism first.} 50k/s over 3 Kafka partitions cannot feed enough Spark tasks; the design scales to 96 partitions (peak headroom 3x) so each micro-batch parallelizes. \texttt{maxOffsetsPerTrigger} caps batch size so a peak burst stretches latency slightly instead of OOMing executors---backpressure by configuration, monitored by the input-vs-processing-rate metric.

\textbf{One table, two consumers.} Bronze is liquid-clustered on \texttt{(event\_hour, campaign\_id)}; pacing dashboards read the last hour (a small, clustered slice); nightly billing scans by hour partition. Micro-batches at 30 seconds with coordinated predictive optimization keep file sizes healthy (Chapter 5's lesson, streaming edition).

\begin{pitfall}
\textbf{Measuring latency at the dashboard only.} End-to-end latency is producer-to-Bronze plus dashboard refresh. Instrument both: \texttt{lastProgress} batch lag for the stream, refresh timestamp for the dashboard. A ``slow pipeline'' was, in the real incident this mirrors, a dashboard refreshing every 5 minutes against a healthy 30-second stream.
\end{pitfall}

\subsection{Risks and Mitigations}
\begin{table}[H]
\centering
\small
\begin{tabular}{@{}p{4.6cm}p{7.6cm}@{}}
\toprule
\textbf{Risk} & \textbf{Mitigation} \\
\midrule
Peak burst triples normal rate & Partition headroom and capped offsets; lag-based autoscaling review \\
Retention gap truncates history & Replay drills sized to retention; \texttt{failOnDataLoss} policy documented \\
Dashboard looks real-time but is stale & Freshness tile on every real-time dashboard \\
Stream failure pages the wrong team & Ownership tags on checkpoints and jobs; alert routing tested \\
Billing reads the same table mid-write & Delta snapshot isolation; documented writer slices \\
\bottomrule
\end{tabular}
\caption{Week 13 scenario risks and mitigations.}
\label{tab:ch13-risks}
\end{table}

\section{Design Decision Guide}
\index{decision guide!Week 13}%
Streaming's first decisions are semantic; the configuration follows from them.

\begin{table}[H]
\centering
\small
\begin{tabular}{@{}p{3.4cm}p{4.4cm}p{4.6cm}@{}}
\toprule
\textbf{Decision} & \textbf{Choose the first when} & \textbf{Choose the second when} \\
\midrule
Append vs.\ update output mode & Event facts, new rows only & Keyed results that evolve \\
\texttt{availableNow} vs.\ processing-time & Minutes-freshness, cost-aware & Continuous low-latency flow \\
Micro-batch vs.\ continuous & Nearly everything & Measured sub-second mandates \\
New checkpoint vs.\ reuse & Query semantics changed & Never reuse across definitions \\
Delta sink vs.\ external sink & Lakehouse consumers & Operational systems needing push \\
\texttt{failOnDataLoss} true vs.\ false & Gaps must halt the pipeline & Retention gaps are tolerable and logged \\
\bottomrule
\end{tabular}
\caption{Week 13 decision guide.}
\label{tab:ch13-decisions}
\end{table}

Write the latency requirement down before choosing a trigger---``real-time'' unquantified becomes the most expensive setting by default. The checkpoint governance rules from the chapter override every convenience.

\section{Operational Guidance for Week 13}
\index{operational guidance!Week 13}%
\subsection{Performance and Reliability}
Alert on the two rates---input versus processing---and on batch-duration trend; a stream that takes 25 seconds per 30-second trigger today falls over at the next traffic step. Back up checkpoint directories with the same care as tables (they are the recovery state), and document the new-checkpoint procedure (query version bump + backfill plan) before you need it. Upgrade runtimes against a copy of the stream first: state-store formats evolve.

\subsection{Security and Governance Checklist}
\begin{itemize}
    \item Checkpoints and schema locations in governed volumes, one per query, named in the runbook.
    \item Source credentials (Kafka, cloud notifications) from secret scopes only.
    \item Sink table writes restricted to the streaming service principal.
    \item \texttt{failOnDataLoss} choice documented per stream with the retention window it assumes.
\end{itemize}

\subsection{Troubleshooting Playbook}
\begin{table}[H]
\centering
\small
\begin{tabular}{@{}p{3.6cm}p{4.2cm}p{4.8cm}@{}}
\toprule
\textbf{Symptom} & \textbf{Likely cause} & \textbf{First action} \\
\midrule
Stream falls behind over hours & Parallelism or batch-size bound & Check rates in lastProgress; partitions, then cap, then executors \\
Restart fails after code change & State schema incompatible with checkpoint & New checkpoint + planned backfill; never force the flag casually \\
Duplicate rows after crash & Sink not idempotent for replays & MERGE/replace sink; verify checkpoint-sink pairing \\
Batch 0 huge, then normal & Backlog on first start & Expected; cap with maxFiles/OffsetsPerTrigger \\
\bottomrule
\end{tabular}
\caption{Week 13 troubleshooting quick reference.}
\label{tab:ch13-trouble}
\end{table}

\section{Interview Q\&A}
\subsection{Concepts and Trade-offs}
\begin{enumerate}
    \item \textbf{Q: Explain micro-batch execution to a backend engineer.}\\
    A: The stream is an infinite table; each trigger runs a small incremental batch query over rows since the last committed offset, writes transactionally, and commits offsets to the checkpoint. Latency is the trigger interval plus batch time.
    \item \textbf{Q: What exactly does the checkpoint make safe?}\\
    A: Restart position (offsets), stateful-operator state, and query metadata. Combined with an idempotent/transactional sink, you get exactly-once; without the sink half you get at-least-once with duplicates.
    \item \textbf{Q: Append versus update output mode?}\\
    A: Append for event facts; update for keyed results that change over time (and then only to sinks that accept updates, or via \texttt{foreachBatch} MERGE into Delta); complete rewrites the whole result---small dimension-style outputs only.
    \item \textbf{Q: When is \texttt{availableNow} the right trigger?}\\
    A: Whenever the freshness requirement is minutes-to-hours: you get streaming sources with batch economics---process the backlog, stop, pay nothing until the next schedule.
    \item \textbf{Q: A stream keeps up off-peak but falls behind at peak. Options?}\\
    A: Check batch duration composition first: add source partitions (parallelism), cap batch size (backpressure), increase executors, or remove an accidental wide transformation. Scale the bottleneck, not the whole pipeline.
    \item \textbf{Q: How does Structured Streaming differ from ``running a batch job every minute''?}\\
    A: The checkpointed offsets, incremental state, and exactly-once composition---a cron batch re-reads, re-processes, and re-writes; the stream processes each record once and remembers it did.
    \item \textbf{Q: What metric would you graph on a stream's dashboard first?}\\
    A: Input rate versus processing rate per batch---it is the earliest, cleanest leading indicator of every streaming failure mode.
    \item \textbf{Q: What is the smallest safe streaming change process?}\\
    A: Copy the query, new checkpoint, replay into a recovery table, reconcile, then cut over. The query-under-test never shares state with the one in production.
\end{enumerate}

\section{Further Reading}
The Structured Streaming Programming Guide is the reference \cite{sparkStructuredStreaming}; the Kafka integration guide previews Chapter 15 \cite{sparkKafkaIntegration}. Delta-as-sink semantics are covered in the Delta documentation \cite{deltaLakeDocs,databricksDelta}, and Auto Loader (Chapter 11) is the file source used throughout \cite{databricksAutoLoader}.

\section*{Key Takeaways}
\begin{itemize}
    \item Structured Streaming is incremental batch over an unbounded table---same API, same optimizer, same sinks.
    \item Exactly-once = checkpoint (offsets + state) plus an idempotent sink; neither half suffices alone.
    \item Output mode follows query semantics; trigger choice is a cost-latency dial.
    \item One checkpoint per query; governed locations; new checkpoint for incompatible changes.
    \item Monitor input rate versus processing rate; fix parallelism and batch size before hardware.
\end{itemize}

\section{Exercises}
\begin{enumerate}
    \item \textbf{(Conceptual)} Why does exactly-once require cooperation between checkpoint and sink?
    \item \textbf{(Conceptual)} Which output mode can a grouped aggregation use with a Delta sink, and why is \texttt{complete} usually wrong at scale?
    \item \textbf{(Hands-on)} Complete the Thursday lab, including the 1500-row uniqueness proof.
    \item \textbf{(Hands-on)} Add \texttt{maxFilesPerTrigger} and observe the latency/throughput trade-off over a synthetic burst.
    \item \textbf{(Hands-on)} Write the listener-based metrics table and chart batch duration over 20 batches.
    \item \textbf{(Design)} Size Kafka partitions and trigger interval for the Friday scenario at 3x peak; show the arithmetic.
    \item \textbf{(Operations)} Write the alert rule for ``processing rate $<$ 80\% of input rate for 10 minutes.''
    \item \textbf{(Architecture)} When would you choose DLT streaming tables (Chapter 10) over this hand-built stream? What changes operationally?
\end{enumerate}
